\chapter{Resultados}
\label{ch:resultados}

Este capítulo apresenta os resultados das duas frentes de coleta descritas no Capítulo~\ref{ch:metodologia}: as entrevistas com gestores dos três estabelecimentos e o questionário com frequentadores. Ao final, as duas fontes são integradas conforme o plano de análise da Seção~\ref{sec:analise}.

\section{Casos estudados}
\label{sec:res-casos}

\pendente{caracterizar os três estabelecimentos (perfil, bairro, ano de abertura, porte, modalidade de autoatendimento) e resumir cada entrevista após a transcrição}

\section{Questionário com frequentadores}
\label{sec:res-questionario}

\subsection{Participação e perfil da amostra}
\label{sec:res-perfil}

\pendente{número total de respostas; respostas excluídas por falta de consentimento; tamanho da base principal (Porto Alegre e região metropolitana) e da base de outras localidades; número de respondentes que nunca encontraram cardápios digitais}

% Tabela sugerida: perfil da base principal (faixa etária, gênero, escolaridade,
% frequência de ida, tipo de estabelecimento, deficiência), com n e %.

\pendente{tabela e comentário sobre o perfil da amostra, comparando com o público esperado e registrando o viés da amostragem por conveniência}

\subsection{Experiência e preferência}
\label{sec:res-experiencia}

\pendente{frequência com que os respondentes encontram cardápios digitais e preferência geral (físico, digital, indiferente)}

\subsection{Construtos UTAUT2 e extensões contextuais}
\label{sec:res-construtos}

% Tabela sugerida: por construto, média, mediana, desvio-padrão, % de
% concordância (4 e 5) e correlação entre os dois itens, após recodificar
% EE2, HM2 e PV2. Extensões (TRU, PRI, SOC) item a item.

\pendente{tabela de escores por construto e comentário sobre os construtos mais e menos favoráveis}

\subsection{Problemas técnicos e situações de uso}
\label{sec:res-problemas}

\pendente{frequência de cada problema técnico relatado, taxa de desistência ou frustração e categorias das respostas abertas (funciona melhor, funciona pior, relatos de frustração)}

\subsection{Cruzamentos exploratórios}
\label{sec:res-cruzamentos}

\pendente{escores e preferência por faixa etária, frequência de ida, deficiência ou dificuldade e localidade (Porto Alegre e região metropolitana versus outras cidades)}

\section{Integração das fontes}
\label{sec:res-integracao}

\pendente{para cada categoria das entrevistas, indicar se o questionário traz evidência convergente, divergente ou complementar}
